\documentclass[12pt, twoside]{article}
\input{../preamble}

\fancyhead[LO]{La Scuola d'Italia / Dr.\ Huson / IB Math \\* 30 March 2026}
\fancyhead[RO]{Markscheme}
\fancyfoot[C]{\thepage}

\begin{document}
\subsubsection*{6.4 Classwork: Systems of Equations --- Markscheme}

\begin{enumerate}[itemsep=0.2cm]

%--- Problem 1: Linear system [7 marks] ---
\item The diagram shows two lines, $y = 2x - 1$ and $y = -x + 5$.

\begin{enumerate}[itemsep=0.3cm]
  \item Write down the gradient of each line. \hfill [2]

  \begin{minipage}[t]{0.62\textwidth}
  Gradient of $y = 2x - 1$: \quad $m = 2$\\[0.15cm]
  Gradient of $y = -x + 5$: \quad $m = -1$
  \end{minipage}%
  \begin{minipage}[t]{0.35\textwidth}\raggedleft
  \textbf{(A1)}\\[0.15cm]
  \textbf{(A1)}
  \end{minipage}

  \medskip
  \item Find the coordinates of the point of intersection by solving
        the system algebraically. \hfill [3]

  \begin{minipage}[t]{0.62\textwidth}
  $2x - 1 = -x + 5$\\[0.15cm]
  $3x = 6 \Rightarrow x = 2$\\[0.15cm]
  $y = 2(2) - 1 = 3$\\[0.15cm]
  Point of intersection: $(2,\; 3)$
  \end{minipage}%
  \begin{minipage}[t]{0.35\textwidth}\raggedleft
  \textbf{(M1)}\\[0.15cm]
  \textbf{(A1)}\\[0.15cm]
  ~\\[0.15cm]
  \textbf{(A1)}
  \end{minipage}

  \medskip
  \item Verify that the point satisfies both equations. \hfill [2]

  \begin{minipage}[t]{0.62\textwidth}
  Check $y = 2x - 1$: $y = 2(2) - 1 = 3$ \quad \checkmark\\[0.15cm]
  Check $y = -x + 5$: $y = -(2) + 5 = 3$ \quad \checkmark
  \end{minipage}%
  \begin{minipage}[t]{0.35\textwidth}\raggedleft
  \textbf{(M1)}\\[0.15cm]
  \textbf{(A1)}
  \end{minipage}
\end{enumerate}

\newpage

%--- Problem 2: Quadratic system [8 marks] ---
\item The diagram shows two parabolas, $y = x^2 - 4x + 3$ and
$y = -x^2 + 2x + 3$.

\begin{enumerate}[itemsep=0.3cm]
  \item Show that the $x$-coordinates of the points of intersection
        satisfy $2x^2 - 6x = 0$. \hfill [2]

  \begin{minipage}[t]{0.62\textwidth}
  $x^2 - 4x + 3 = -x^2 + 2x + 3$\\[0.15cm]
  $2x^2 - 6x = 0$
  \end{minipage}%
  \begin{minipage}[t]{0.35\textwidth}\raggedleft
  \textbf{(M1)}\\[0.15cm]
  \textbf{(A1)} \textbf{(AG)}
  \end{minipage}

  \medskip
  \item Hence find the $x$-coordinates of the points of intersection. \hfill [2]

  \begin{minipage}[t]{0.62\textwidth}
  $2x(x - 3) = 0$\\[0.15cm]
  $x = 0$ or $x = 3$
  \end{minipage}%
  \begin{minipage}[t]{0.35\textwidth}\raggedleft
  \textbf{(A1)}\\[0.15cm]
  \textbf{(A1)}
  \end{minipage}

  \medskip
  \item Find the coordinates of both points of intersection. \hfill [2]

  \begin{minipage}[t]{0.62\textwidth}
  $x = 0$: \quad $y = 0 - 0 + 3 = 3$ \quad $\Rightarrow (0,\; 3)$\\[0.15cm]
  $x = 3$: \quad $y = 9 - 12 + 3 = 0$ \quad $\Rightarrow (3,\; 0)$
  \end{minipage}%
  \begin{minipage}[t]{0.35\textwidth}\raggedleft
  \textbf{(A1)}\\[0.15cm]
  \textbf{(A1)}
  \end{minipage}

  \medskip
  \item Find the coordinates of the vertex of $y = x^2 - 4x + 3$. \hfill [2]

  \begin{minipage}[t]{0.62\textwidth}
  $x = \dfrac{-(-4)}{2(1)} = 2$\\[0.15cm]
  $y = 4 - 8 + 3 = -1$\\[0.15cm]
  Vertex $= (2,\; -1)$
  \end{minipage}%
  \begin{minipage}[t]{0.35\textwidth}\raggedleft
  \textbf{(M1)}\\[0.15cm]
  ~\\[0.15cm]
  \textbf{(A1)}
  \end{minipage}
\end{enumerate}

\end{enumerate}
\end{document}
